%\section{Multi-party NIKE in the Bounded Storage Model with Background and (Full)Security}
%In this section, we give the definitions and theorems used to construct our multi-party NIKE in the bounded storage model, along with the proof of its security.
%\subsection{Background of Bounded Storage Model}
%\section{Proof of \cref{lem:pairwise}}
%\label{sec:pairwise}

\begin{proof}
%Let $ t_i^j=h_i(j)$ for all $i\in[\pnum]$ and $j\in\Z_q$ and $( u_j)_{j\in\Z_q}$ be the sequence such that $ u_j= h_1(j)$ if there exist indices $j_2,\cdots,j_{\pnum}$ satisfying
%$\msb_\kk( h_1(j))=\msb_\kk( t_i^{j_2})=\cdots=\msb_\kk( t_\pnum^{j_\pnum})$ and $ u_i=\bot$ otherwise.
Let $\hx_1,\hx_2$ be any distinct elements in $GF(n)$.
%To prove this lemma, it is sufficient to show that 
%the sequence $( u_i)_{i\in[q]}$ restricted to $ u_i\neq\bot$ is $\kk$-approximately pairwise independent, i.e. $( u_i)_{i\in[q]}$ satisfies
%$
%\Pr[ u_i=\hx_1\land  u_j=\hx_2| u_i\neq \bot\land  u_j\neq\bot]=\frac{1}{k(k-1)}
%$
%for all $i,j\in\Z_q$. 
Due to approximately pairwise independence, we have 
$$
\Pr[ h_\pnum(i)=\hx_1\land  h_\pnum(j)=\hx_2]=1/(n(n-1))
$$
for all distinct $i,j\in[q]$,
and
$$
\Pr[\exists i_l,j_l: h_l(i_l)=\hx_1\land h_l(j_l)=\hx_2]=q(q-1)/(n(n-1))
$$
for all $l\in[\pnum-1]$.
Hence, for any distinct $i,j\in\Z_q$, we have
\[\begin{split}
\Pr[ u_i=\hx_1\land  u_j=\hx_2]
=&\Pr[\exists (i_l,j_l)_{l\in[\pnum-1]}: 
 h_1(i_1)=\cdots= h_{\pnum-1}(i_{\pnum-1})=\hx_1= h_\pnum(i)\land\\
 &\tab h_1(j_1)=\cdots= h_{\pnum-1}(j_{\pnum-1})=\hx_2= h_\pnum(j)]\\
=&\Pr[ h_\pnum(i)=\hx_1\land  h_\pnum(j)=\hx_2]\cdot\prod_{l=1}^{\pnum-1}\Pr[\exists i_l,j_l: h_l(i_l)=\hx_1\land  h_l(j_l)=\hx_2]\\
%=\Pr[&t_1^i=x_1\land h_1(j)=x_2]\cdot\Pr[\exists i_2,\cdots,i_\pnum,j_2,\cdots,j_\pnum: \\&
%\msb_\kk(t_2^{i_2})=\cdots=\msb_\kk(t_\pnum^{i_\pnum})=\msb_\kk(x_1)\land\\
%&\msb_\kk(t_2^{j_2})=\cdots=\msb_\kk(t_\pnum^{j_\pnum})=\msb_\kk(x_2)]\\
=&1/(n(n-1))\cdot \Big(q(q-1)/(n(n-1))\Big)^{\pnum-1},
\end{split}\]
and
\[\begin{split}
\Pr[ u_i\neq\bot\land  u_j\neq \bot]
=&\Pr[\exists (i_l,j_l)_{l\in[\pnum-1]}: h_1(i_1)=\cdots= h_{\pnum-1}(i_{\pnum-1})= h_\pnum(i)\land\\
&\tab  h_1(j_1)=\cdots= h_{\pnum-1}(j_{\pnum-1})= h_\pnum(j)]
\\
=&\prod_{l=1}^\pnum\Pr[\exists i_l,j_l: h_l(i_l)=h_\pnum(i)\land  h_l(j_l)= h_\pnum(j)]
=\Big(q(q-1)/(n(n-1))\Big)^{\pnum-1},
\end{split}\]
i.e.,
$$
 \Pr[ u_i=\hx_1\land  u_j=\hx_2\mid u_i\neq \bot\land  u_j\neq\bot]=\frac{\Pr[ u_i=\hx_1\land  u_j=\hx_2]}{\Pr[ u_i\neq\bot\land  u_j\neq \bot]}=1/(n(n-1)).
$$
Therefore, conditioned on two surviving positions, their values are uniformly distributed over ordered pairs of distinct field elements. Restricting to the non-$\bot$ entries and then taking the first $\lambda$ entries preserves this pairwise distribution. Hence, conditioned on $\mathsf{Good}$, $\vec{s}_{\lambda}$ satisfies the condition required by \cref{thm:cm}.
\end{proof}


%In this section, we give the backgrounds used to construct our multi-party NIKE in the bounded storage model.
%\heading{Pairwise independence.}
%\label{sec:pairwisedef}
%In this section, we recall the definitions of uniformly and approximately pairwise independence, (approximately) strongly 2-universal hash function family, and 2-universal hash function family, and recall their instantiations in \cref{fig:con-hash}.~\footnote{The approximate pairwise independence and approximate 2-universal hash function family are defined in \cite{FOCS:Zuckerman91,C:CacMau97} in an implicit way.}
%%We refer the reader to~\cref{sec:pairwise} for the definitions of uniformly and approximately pairwise independence, (approximately) strongly 2-universal hash function family, and 2-universal hash function family, and recall their instantiations in \cref{fig:con-hash}.
%%Roughly approximately pairwise independent variables are uniformly pairwise independent conditioned on that they are all distinct.
%
%\begin{definition}[Uniform pairwise independence]
%\label{def:upairwise}
%A sequence of random variables $x_1,\cdots,x_q$ with alphabet $\mathcal Y$ are uniformly pairwise independent if for any $1\leq i<j\leq q$ and any $\hat y_1, \hat y_2 \in \mathcal X$, we have
%	$\Pr[x_i = \hat y_1 \land x_j = \hat y_2] = 1/|\mathcal{Y}|^2$.
%\end{definition}
%
%
%\begin{definition}[Approximate pairwise independence \cite{FOCS:Zuckerman91,C:CacMau97}]
%\label{def:apairwise}
%A sequence of random variables $x_1,\cdots,x_q$ with alphabet $\mathcal Y$ are approximate pairwise independent if for any $1\leq i<j\leq q$ and any distinct $\hat y_1, \hat y_2 \in \mathcal Y$, we have
%$\Pr[x_i = \hat y_1 \land x_j = \hat y_2] = 1/|\mathcal{Y}|(|\mathcal{Y}|-1)$.
%%The definition of \emph{approximate pairwise independence} is exactly the same as \cref{def:upairwise} except that we replace 
%%``$\frac{1}{|\mathcal{Y}|^2}$''
%%by
%%``$\frac{1}{|\mathcal{Y}|(|\mathcal{Y}|-1)}$''.
%\end{definition}
%
%
%%For any $q\in\N$, a sequence of random variables $\vec x_1,\cdots,\vec x_q$ from a range $\mathcal X$ are pairwise independent if for any $1\leq i<j\leq q$ and any $\hy_1, \hy_2 \in \bin^\nn$ with $\hy_1\neq \hy_2$, we have
%%	$$\Pr[\vec x_i = \hy_1 \land \vec x_j = \hy_2] = \frac{1}{|\mathcal{X}|^2}.$$
%
%\begin{definition}[Strongly 2-universal hash function family]
%\label{def:shash}
%	A set $\mathcal H$ of functions with domain $\mathcal X$ and range $\mathcal Y$ is said to be a \emph{strongly 2-universal hash function family} if for any distinct $x_1,x_2 \in \mathcal{X}$, $h(x_1)$ and $h(x_2)$ are uniformly pairwise independent for $h\getsr \mathcal H$.
%\end{definition}
%
%\begin{definition}[Approximate strongly 2-universal hash function \cite{FOCS:Zuckerman91,C:CacMau97}]
%\label{def:ahash}
%The definition of \emph{approximate strongly 2-universal hash function family} is exactly the same as \cref{def:shash} except that we replace ``uniformly pairwise independent'' by ``approximate pairwise independent''.
%\end{definition}
%
%
%\begin{definition}[2-universal hash function family]
%\label{def:hash}
%	A set $\mathcal H$ of functions with domain $\mathcal X$ and range $\mathcal Y$ is said to be a \emph{2-universal hash function family} if for any distinct $x_1,x_2 \in \mathcal{X}$, we have
%	$\Pr[g(x_1)=g(x_2)]\leq 1/|\mathcal Y|$.
%\end{definition}

%We now recall a main theorem in \cite{C:CacMau97} for privacy amplification. 
